\documentclass{situationist-preprint}

\title{Socrates and the Duty to Remain}
\msid{ETH-17-b}
\msfamily{ETH-17}
\msdiscipline{Normative and Applied Ethics}
\mstype{Companion manuscript}
\mskeywords{Socrates; Crito; civil disobedience; acceptance of penalty; unjust verdicts; fair play}

\begin{document}
\maketitle

\begin{abstract}
Must one submit to an unjust verdict reached by a just procedure? Socrates in the \emph{Crito} held that he must, and the argument of the Laws is often read as asserting a strong general duty of submission. We argue that there is no such general duty: the duty to comply with the outcomes of a just procedure is the weak, fairness-based duty to obey, and it can be outweighed. Socrates nonetheless had a duty to remain, of a more specific kind. He had openly refused to abandon conduct the city condemned, for the sake of his fellow citizens, and the acceptance condition of ETH-17-a binds those who breach in this way to submit to the penalty the process imposes, even when they judge the verdict mistaken.
\end{abstract}

\section{Two readings of the Crito}

In the \emph{Crito} the Laws argue that Socrates, having lived in Athens and enjoyed its benefits, agreed to abide by its judgments, and that by escaping he would be trying, so far as in him lies, to destroy them \citep[50a--54d]{PlatoCrito}. Read as a general thesis, this asserts a strong duty to submit to any verdict reached by a just procedure, however unjust its content. That reading sits uneasily with the \emph{Apology}, where Socrates tells the jury that he would not abandon philosophy even if acquittal were made conditional on it, since he must obey the god rather than them \citep[29c--d]{PlatoApology}. The tension is much discussed \citep{Woozley1979,Kraut1984}; ETH-17-a gives a resolution.

\section{The argument}

ETH-17-a grounds the duty to obey in fairness (its Assumption~2.2, the conclusion of POL-03). Submission to verdicts is one of the burdens of a cooperative scheme, but the duty so grounded is weak and defeasible, and nothing in its ground makes it stronger when its object is a verdict rather than a statute.

\begin{proposition}\label{p:nogeneral}
There is no general duty to submit to an unjust verdict reached by a just procedure that is stronger than the weak duty to obey just laws.
\end{proposition}

\begin{proof}
A duty to submit to verdicts is either grounded in fairness or not. If it is, its strength is that of Assumption~2.2 of ETH-17-a, which is weak and defeasible. If it is grounded otherwise, in consent or in gratitude as the Laws also suggest, it inherits the defects of those grounds as accounts of political obligation \citep{Simmons1979}, which POL-03 rejects. Either way there is no stronger duty.
\end{proof}

Proposition~\ref{p:nogeneral} does not by itself license escape: whether a grave harm to an innocent person outweighs the weak duty depends on the case, and the four conditions of ETH-17-a (its Definition~2.3), which concern harm to others, do not settle it. Socrates' case has a further feature that does.

\begin{proposition}\label{p:remain}
An agent who has openly breached, or openly refused to comply with, a demand of a just scheme, for the sake of others, has a duty to submit to the penalty the scheme's just procedure imposes for that conduct, even if she judges the verdict mistaken.
\end{proposition}

\begin{proof}
By Principle~P1 of ETH-17-a, a breach of a just law is unfair if the agent withdraws from the burden that the scheme attaches to its breach. The burden the scheme attaches is the penalty imposed by its procedure, not the penalty the agent judges correct; if each disobedient could substitute her own verdict, acceptance would be acceptance of nothing. So an agent who breaches openly and then evades the penalty free-rides, and her breach fails the acceptance condition (Definition~2.3(iv) of ETH-17-a). Since that condition is part of what makes the breach permissible (Theorem~3.3 of ETH-17-a), she has reason to submit; and the reason is a duty owed to those who bear the scheme's burdens.
\end{proof}

Socrates had declared before the court that he would continue to examine his fellow citizens whatever the verdict, and he presented this as a service to them \citep[30a--31c]{PlatoApology}. His conduct was open defiance for the sake of others, and Proposition~\ref{p:remain} applies to it. His duty to remain was the acceptance condition, not a general duty of submission. This is why the Laws' argument persuades in his case and fails as a general thesis.

\begin{objection}
The verdict was unjust, and Socrates' execution was a grave harm. If the gravity of harm can justify breaching a law, it can justify escaping a sentence.
\end{objection}

\begin{reply}
Escape would be a second breach whose aim is to evade the penalty for the first; it fails the acceptance condition by its nature, and makes the first breach a case of free-riding. Proposition~\ref{p:remain} applies, however, only where the procedure is just; a verdict reached by a sham procedure imposes no burden of a just scheme, and the duty to remain lapses.
\end{reply}

\section{Conclusion}

There is no general duty to submit to unjust verdicts beyond the weak duty to obey. The disobedient who acts openly for others' sake has a specific duty to remain, the acceptance condition applied to its own consequences, and Socrates was right to stay for that reason.

\begin{thebibliography}{9}\small
\bibitem[Kraut(1984)]{Kraut1984} Kraut, R. (1984). \emph{Socrates and the State}. Princeton University Press.
\bibitem[Plato(Apology)]{PlatoApology} Plato. \emph{Apology}. Cited by Stephanus page.
\bibitem[Plato(Crito)]{PlatoCrito} Plato. \emph{Crito}. Cited by Stephanus page.
\bibitem[Simmons(1979)]{Simmons1979} Simmons, A.~J. (1979). \emph{Moral Principles and Political Obligations}. Princeton University Press.
\bibitem[Woozley(1979)]{Woozley1979} Woozley, A.~D. (1979). \emph{Law and Obedience: The Arguments of Plato's Crito}. London: Duckworth.
\end{thebibliography}

\end{document}
